% 2010 年安徽卷理科第 8 题：几何体的三视图（正视图、侧视图、俯视图）。
% 来源 PDF：普通高考/2010/2010安徽理.pdf，SHA-256 2760eb9e09374139a51303b0c2590746c74dfef7fed99ac6b2eace305af646c5。
% 原页：第 1 页（A3 三栏）第 2 栏上部，书页码 569。
% 重绘依据：下部长方体 8×10×2；上部长方体 6×2×8，居中置于下部之上。
% 正视图宽 8、高 2，上部宽 6、两侧各 1；侧视图深 10、高 2，上部深 2、偏右 2；
% 俯视图 8×10，内框 6×2、距下边 2。
\begin{tikzpicture}[semithick, scale=0.38, every node/.style={inner sep=1.5pt, font=\small}]
  % 正视图
  \draw (0,0) rectangle (8,2);
  \draw (1,2) rectangle (7,10);
  \draw[<->] (0,10.45) -- node[above] {$1$} (1,10.45);
  \draw[<->] (1,10.45) -- node[above] {$6$} (7,10.45);
  \draw[<->] (7,10.45) -- node[above] {$1$} (8,10.45);
  \draw[<->] (-0.45,0) -- node[left] {$2$} (-0.45,2);
  \draw[<->] (-0.45,2) -- node[left] {$8$} (-0.45,10);

  % 侧视图
  \draw (11,0) rectangle (21,2);
  \draw (17,2) rectangle (19,10);
  \draw[<->] (11,10.45) -- node[above] {$6$} (17,10.45);
  \draw[<->] (17,10.45) -- node[above] {$2$} (19,10.45);
  \draw[<->] (19,10.45) -- node[above] {$2$} (21,10.45);
  \draw[<->] (21.45,0) -- node[right] {$2$} (21.45,2);
  \draw[<->] (21.45,2) -- node[right] {$8$} (21.45,10);

  % 俯视图
  \draw (4,-14) rectangle (12,-4);
  \draw (5,-12) rectangle (11,-10);
  \draw[<->] (12.45,-4) -- node[right] {$6$} (12.45,-10);
  \draw[<->] (12.45,-10) -- node[right] {$2$} (12.45,-12);
  \draw[<->] (12.45,-12) -- node[right] {$2$} (12.45,-14);
  \draw[<->] (4,-14.45) -- node[below] {$8$} (12,-14.45);

  \node at (4,-2.25) {正（主）视图};
  \node at (16,-2.25) {侧（左）视图};
  \node at (8,-16.2) {俯视图};
\end{tikzpicture}
